\chapter{小行星带的诞生}

投送筒消失后，残舟里安静得反常。

警报还在响，推进环还在咳出断续火光，舱壁也仍然发热。可洛安觉得那些声音全都退到很远的地方。怀里空了以后，他反倒不知道该把手放在哪里。

岚柯靠在舱壁上，闭着眼，像终于允许自己虚弱一下。索勒守在控制台前，手指还搭在姿态调节杆上，可他已经没有多少可调的东西了。

“它过去了吗？”洛安问。

索勒看着屏幕上最后一点蓝光残影。

“过去了。”

没人欢呼。

这句话只是在确认一件已经付出太多代价的事。

残舟开始偏离窗口边缘，慢慢滑向法厄同外轨道的碎片流。洛安没有再挣扎。他转头看向舷窗，看见自己的世界正在裂成一条巨大的环。

法厄同已经失去球体形状。

远处大陆像烧碎的甲片一样从核心周围剥离，城市、山脉、海盆、轨道塔和地下设施全被扯进不同方向。某些碎片仍带着灯光，某些碎片拖着大气残痕，像从伤口里飘出来的雾。更大的岩块互相碰撞，撞击后碎成更小的星尘。

这不是爆炸后的瞬间混乱。

它正在形成秩序。

查理在洛安的视野里看着这一切，心底升起一种难以形容的寒意。碎片并没有完全飞散。木星的引力、太阳的牵引、法厄同残余速度，把它们慢慢拉成一条宽阔的带。那条带一开始粗糙、暴烈、充满撞击，可在更漫长的时间里，它会被磨平，变冷，变成一个看上去仿佛天然存在的结构。

小行星带。

这个词太轻了。

它只有四个字，干净、冷静，适合写进教材和星图。可此刻在舷窗外，每一块碎片都带着某个场所的残余重量。那里可能有一条震翼族航路的转接塔，有鳞裔族尚未转移完的库房，有节肢族断开的节点，也有渊民族早已干涸的水厅。它们未来会被编号、测量、归类，成为冰、镍铁、硅酸盐和轨道偏心率。

没有人会在参数里听见尖叫。

这个词从查理的记忆里浮起时，洛安也感到了。

“你知道它会变成什么？”洛安问。

查理沉默了一下。

“知道一点。”

“说。”

“很久以后，那里会只剩碎石。我的时代会把它叫作小行星带。很多人会研究它的轨道、成分和碰撞历史。也许有人会猜这里曾经有一颗星球，但大多数人只会把它当成太阳系结构的一部分。”

洛安看着舷窗外那些仍亮着灯的碎片。

“他们不会知道我们在这里？”

“不会。”

这两个字说出来，比查理预想中更疼。

洛安没有骂，也没有笑。他只是很慢地点了一下头。

“也好。”

“什么也好？”

“如果他们不知道，就不会把这里写成谁赢了。”洛安说，“也不会有人拿我们的尸体继续讲自己的荣耀。”

查理说不出话。

残舟被一股引力波扫过，舱体发出长长的撕裂声。岚柯睁开眼，索勒重新握紧控制杆。可洛安仍看着那条正在形成的碎片带。

他看见法厄同的天空变成石头。

看见战争的旗帜、律令、巢群、地底库、水声和震翼族航线，全都被卷进同一片沉默。

也许这就是巨鲸留下的另一种记录。

没有碑文，没有胜利者，没有一族能把整颗星球的死亡占为己有。

只有碎片绕着太阳慢慢转。

转到所有名字都失去发音。

转到后来者抬头看见它们时，只觉得那里本来就该有一条带。

洛安看了很久，终于低声说：“如果未来有人真的猜到这里曾经有星球，别让他们只猜到石头。”

查理回答：“我会记住。”

他不知道自己醒来后还能记住多少。可这一刻，他确实把那条正在形成的碎片带刻进了意识深处，像替一个没有墓碑的世界，记住第一行墓志铭。
